%% THERMODYNAMIQUE %%
%
\section{Rapport de m\'elange en vapeur d'eau \`a partir de l'humidit\'e
relative}

\itemstyle{ Op\'erateur~: } {\sf RapM\_Hu}
\index{RapM\_Hu}

%%%%%OBJET%%%%%
\itemstyle{ Objet~: }
Calcul du rapport de m\'elange en vapeur d'eau \`a partir
de l'humidit\'e relative de l'air et des variables
d'�tat (P, T).

Comme pour la fonction sym�trique ({\sf Hu\_RapM}), le calcul de la
tension de vapeur saturante ($ew(T)$) se fait TOUJOURS par 
rapport � de l'eau liquide.

%%%%%ARGUMENTS%%%%%
\bigskip
%%%%%ENTREES%%%%
\begin{tabarg}
\multicolumn{3}{l}{\titretab{Entr\'ees}}\\
\hline
$H_u$ & Vecteur, {\tt float} & Humidit\'e relative ($\%$)\\
$T_s$ & Vecteur, {\tt float} & Temp\'erature statique ($K$ ou $^\circ C$) \\
$P_s$ & Vecteur, {\tt float} & Pression statique ($hPa$) \\
\end{tabarg}

\medskip
%%%%%SORTIES%%%%
\begin{tabarg}
\multicolumn{3}{l}{\titretab{Sorties}}\\
\hline
$RapM$ & Vecteur, {\tt float} & Rapport de m\'elange en vapeur d'eau ($g/kg$)\\
\end{tabarg}


%%%%%FORMULE%%%%
\itemstyle{ Formule }

%%% THERMODYNAMIQUE
\[
e = \frac{e_w(T_s) . H_u}{100} 
\]
(formule identique m�me si $T_s < 0$ )
\[
RapM = \frac{R_a}{R_v}.
       \frac{e} 
            {P_s - e }
\]
avec $\frac{R_a}{R_v} = 0.622$ 


%%%%COMMENTAIRE%%%%
\itemstyle{Commentaire~:}

Les fonctions $ew()$ et $ewi()$ sont d\'ecrites en annexe
\ref{annexe-thermo}.

%%%%%AUTEUR%%%%%
\itemstyle{ Auteur~:}
CNRM/GMEI/TRAMM

%%%%%REFERENCES%%%%%
\itemstyle{ R\'ef\'erences~:}
Triplet-Roche\cite{Triplet-Roche}\\
D�finitions des humidit�s (relatives, absolues, sp�cifiques)
et du rapport de m�lange 
selon la table 4.3 des tables m�t�orologiques internationales
de l'OMM \cite{OMM-188.TP.94}.
